\begin{tikzpicture}[
  font=\small,
  >={Stealth[length=2.2mm]},
  error/.style={draw=red!55!black,rounded corners=2pt,fill=red!4,
    minimum height=17mm,text width=34mm,align=center,inner sep=2.5pt},
  defense/.style={draw=blue!60!black,rounded corners=2pt,fill=blue!7,
    minimum height=25mm,text width=34mm,align=center,inner sep=2.5pt},
  validate/.style={draw=orange!75!black,rounded corners=2pt,fill=orange!10,
    minimum height=12mm,text width=118mm,align=center,inner sep=2.5pt},
  errarrow/.style={->,thick,red!65!black},
  block/.style={->,thick,blue!70!black},
  toval/.style={->,densely dashed,orange!80!black}
]
  \node[error] (e1) {\textbf{工位停靠残差}\\位置残差：厘米量级\\航向残差：度量级};
  \node[error,right=8mm of e1] (e2) {\textbf{精细阶段显现}\\位置残差近似等幅传递\\航向项随目标距离增加};
  \node[error,right=8mm of e2] (e3) {\textbf{错误动作持续}\\块内不接收新观测\\可能导致失败或损坏};
  \draw[errarrow] (e1) -- (e2);
  \draw[errarrow] (e2) -- (e3);

  \node[defense,below=10mm of e1] (d1) {\textbf{停靠残差条件化}\\\textbf{模仿学习方法}\\残差状态增广\\随机停靠示范};
  \node[defense,below=10mm of e2] (d2) {\textbf{阶段自适应}\\\textbf{动作块控制方法}\\执行时域调节\\RTC连续衔接};
  \node[defense,below=10mm of e3] (d3) {\textbf{预测一致性驱动的}\\\textbf{运行时容错方法}\\跨块残差检测\\模板恢复或停机};
  \draw[block] (d1) -- node[right,font=\footnotesize]{条件化适应} (e1);
  \draw[block] (d2) -- node[right,font=\footnotesize]{缩短开环} (e2);
  \draw[block] (d3) -- node[right,font=\footnotesize]{截断} (e3);

  \node[validate,below=8mm of d2] (val) {\textbf{整体验证}：C0基线 / C1条件化学习 / C2自适应控制 / C3检测与恢复，相同残差档位与五子任务链，四台机器人重复};
  \draw[toval] (d1.south) -- (d1.south |- val.north);
  \draw[toval] (d2) -- (val);
  \draw[toval] (d3.south) -- (d3.south |- val.north);
  \node[draw=red!50!black,densely dashed,rounded corners=3pt,fit=(e1)(e3),inner sep=2.5mm,
    label={[font=\footnotesize,text=red!60!black]above left:误差传递链}] {};
\end{tikzpicture}
